\documentclass[12pt,a4paper]{article}

\usepackage[T2A]{fontenc}
\usepackage[utf8]{inputenc}
\usepackage[russian]{babel}
\usepackage{amsmath,amssymb,amsfonts}
\usepackage{array}
\usepackage{booktabs}
\usepackage{geometry}
\usepackage{longtable}
\usepackage{pdflscape}
\usepackage{enumitem}
\usepackage{hyperref}

\geometry{margin=2.0cm}
\setlength{\LTpre}{0pt}
\setlength{\LTpost}{0pt}

\newcommand{\sgn}{\operatorname{sign}}
\newcommand{\clip}{\operatorname{clip}}
\newcommand{\softmax}{\operatorname{softmax}}
\newcommand{\Tnorm}[3]{\tanh\!\left(\frac{#1-#2}{#3}\right)}
\newcommand{\spow}[2]{\sgn(#1)\lvert #1 \rvert^{#2}}

\title{Полный свод формул ИИС: V3, V4, V5, V6}
\author{Проект анализа интегрального индекса состояния}
\date{\today}

\begin{document}
\maketitle

\section{Назначение документа}
Этот документ сводит в один файл четыре версии модели ИИС:
\[
V3,\ V4,\ V5,\ V6.
\]
Формулы разделены на три слоя:
\begin{enumerate}[leftmargin=1.2cm]
    \item \textbf{Теоретические} --- структура модели без привязки к конкретным числам калибровки.
    \item \textbf{Практические} --- точные формулы, реально использующиеся в коде для прямых данных.
    \item \textbf{Практико-ориентированные (proxy)} --- формулы, которые включаются при отсутствии части прямых измерений.
\end{enumerate}

\section{Общие обозначения}
\[
\sigma(x)=\frac{1}{1+e^{-x}}, \qquad
\Gamma_+=\max(\Gamma,0), \qquad
V_+=\max(V,0).
\]

\[
\spow{x}{p}=\sgn(x)\lvert x \rvert^p,
\qquad
\softmax(z_1,\dots,z_n)_i=\frac{e^{z_i}}{\sum_{j=1}^{n}e^{z_j}}.
\]

\[
\Tnorm{x}{c}{w}=\tanh\!\left(\frac{x-c}{w}\right).
\]

Базовые физиологические признаки:
\begin{itemize}[leftmargin=1.2cm]
    \item \(P_L, P_R\) --- суммарная мощность левого и правого EEG-блоков;
    \item \(\alpha_{\mathrm{asym}}\) --- alpha-asymmetry;
    \item \(P_\gamma\) --- gamma power;
    \item \(\gamma/\alpha\) --- gamma/alpha ratio;
    \item \(HR\) --- heart rate;
    \item \(HF/LF\) --- отношение компонент HRV;
    \item \(DA\) --- дофаминовый вклад или его proxy;
    \item \(C\) --- кортизоловый вклад или его proxy.
\end{itemize}

\section{Теоретические модели}
\begin{landscape}
\small
\begin{longtable}{>{\raggedright\arraybackslash}p{1.6cm} >{\raggedright\arraybackslash}p{16.7cm}}
\toprule
\textbf{Версия} & \textbf{Теоретическая формула} \\
\midrule
\endfirsthead
\toprule
\textbf{Версия} & \textbf{Теоретическая формула} \\
\midrule
\endhead

\textbf{V3} &
\[
A=\frac{N_L-N_R}{N_L+N_R+\varepsilon},
\qquad
\Gamma=\ln(1+P_\gamma \cdot 10^{12}),
\qquad
V=\ln\!\left(1+\frac{HF}{LF}\right).
\]
\[
Q_{\mathrm{raw}}=\frac{\alpha_L-\alpha_R}{\alpha_L+\alpha_R+\varepsilon},
\qquad
Q=\tanh(\beta_q Q_{\mathrm{raw}}).
\]
\[
K=\frac{R\cdot G}{1+\alpha S^2+\beta F(S)},
\qquad
F(S)=\frac{1}{1+e^{-k_f(S-S_0)}}.
\]
\[
H=\frac{K\cdot DA}{50}-\frac{C}{15}.
\]
\[
IIS_3=\sigma\!\left(0.35A+0.25\Gamma+0.30H+0.10V\right).
\]
Смысл версии: сохранить старое ядро ИИС, но сделать гормональный блок зависящим от коэффициента проницаемости \(K\). \\

\textbf{V4} &
\[
A=\tanh\!\left(
1.10\,\Tnorm{\alpha_{\mathrm{asym}}}{c_\alpha}{w_\alpha}
0.40\,\Tnorm{\dfrac{P_L-P_R}{P_L+P_R+\varepsilon}}{c_A}{w_A}
\right).
\]
\[
\Gamma=\tanh\!\left(
0.30\,\Tnorm{\ln(1+P_\gamma\cdot 10^{12})}{c_\gamma}{w_\gamma}
0.70\,\Tnorm{\ln(1+\gamma/\alpha)}{c_{\gamma/\alpha}}{w_{\gamma/\alpha}}
\right).
\]
\[
V=\tanh\!\left(
0.80\,\Tnorm{HR_0-HR}{0}{w_{hr}}
0.20\,\Tnorm{\ln(1+HF/LF)}{c_{v}}{w_{v}}
\right).
\]
\[
Q_4=\tanh\!\left(0.55A+0.45V-0.25\Gamma\right).
\]
\[
IIS_4=\sigma\!\left(0.10A-0.05\Gamma+0.25V+0.60Q_4\right).
\]
Смысл версии: полностью убрать гормоны и оставить только реально измеряемое EEG/ECG-ядро. \\

\textbf{V5} &
\[
\widehat{A}=\spow{A}{p_q}, \qquad
\widehat{V}=\spow{V}{p_q}, \qquad
\widehat{\Gamma}=\spow{\Gamma}{p_\gamma}.
\]
\[
Q_{\mathrm{lin}}=a_1A+a_2V-a_3\Gamma,
\qquad
Q_{\mathrm{coh}}=AV-b_1|A-V|-b_2\Gamma_+.
\]
\[
S_{AV}=\sgn(\widehat{A}+\widehat{V})\sqrt{|\widehat{A}\widehat{V}|+\varepsilon},
\qquad
E_{AV}=\frac{|\widehat{A}|+|\widehat{V}|}{2}.
\]
\[
C_{AV}=|\widehat{A}-\widehat{V}|+b_3\widehat{\Gamma}_+(1-\widehat{V}_+).
\]
\[
Q_5=\tanh\!\left(
\lambda_1 Q_{\mathrm{lin}}+
\lambda_2 Q_{\mathrm{coh}}+
\lambda_3 S_{AV}+
\lambda_4 E_{AV}\sgn(Q_{\mathrm{lin}})-
\lambda_5 C_{AV}
\right).
\]
\[
\widetilde{A}=\spow{A}{p_i},\qquad
\widetilde{V}=\spow{V}{p_i},\qquad
\widetilde{Q}=\spow{Q_5}{p_i}.
\]
\[
z_5=
r_1\widetilde{A}
r_2\widetilde{V}
r_3\widetilde{Q}
r_4\Gamma-c_0.
\]
\[
S_{\mathrm{reg}}=
\sgn(\widetilde{A}+\widetilde{V})\sqrt{|\widetilde{A}\widetilde{V}|+\varepsilon}
\eta\,
\sgn(\widetilde{Q}+\widetilde{V})\sqrt{|\widetilde{Q}\widetilde{V}|+\varepsilon}.
\]
\[
E_{\mathrm{phase}}=\sqrt{\frac{\widetilde{A}^2+\widetilde{V}^2+\widetilde{Q}^2}{3}},
\qquad
C_{\mathrm{phase}}=|\widetilde{A}-\widetilde{V}|+\xi|\widetilde{Q}-\widetilde{V}|.
\]
\[
contrast_5=
g_1 z_5
g_2 z_5^3
g_3(Q_5-V)
g_4 S_{\mathrm{reg}}
g_5 E_{\mathrm{phase}}\sgn(z_5)
g_6 \Gamma_+
g_7 C_{\mathrm{phase}}
g_8 \Gamma_+(1-V_+).
\]
\[
IIS_{5,\mathrm{base}}=\sigma\!\left(core_5+\mu_1 S_{\mathrm{reg}}-\mu_2 C_{\mathrm{phase}}\right),
\qquad
IIS_{5,\mathrm{contrast}}=0.5+0.5\tanh(contrast_5).
\]
\[
IIS_5=\clip\!\left((1-m)IIS_{5,\mathrm{base}}+mIIS_{5,\mathrm{contrast}},0,1\right).
\]
Смысл версии: сделать шкалу менее плоской и раскрыть режимную геометрию без гормонов. \\

\textbf{V6} &
\[
\widehat{A}=\spow{A}{p}, \qquad
\widehat{V}=\spow{V}{p}, \qquad
\widehat{\Gamma}=\spow{\Gamma}{p_\gamma}, \qquad
\widehat{Q}=\spow{Q_6}{p}.
\]
\[
z^{(Q)}_{\mathrm{reg}}=u_1\widehat{A}+u_2\widehat{V}-u_3\widehat{\Gamma}_+,
\]
\[
z^{(Q)}_{\mathrm{mob}}=u_4\widehat{A}+u_5\widehat{V}-u_6\widehat{\Gamma}_+,
\]
\[
z^{(Q)}_{\mathrm{dep}}=-u_7\widehat{A}-u_8\widehat{V}+u_9\widehat{\Gamma}_+.
\]
\[
\bigl(g^{(Q)}_{\mathrm{reg}},g^{(Q)}_{\mathrm{mob}},g^{(Q)}_{\mathrm{dep}}\bigr)
=\softmax\!\left(\tau_Q z^{(Q)}_{\mathrm{reg}},\tau_Q z^{(Q)}_{\mathrm{mob}},\tau_Q z^{(Q)}_{\mathrm{dep}}\right).
\]
\[
Q_6=
g^{(Q)}_{\mathrm{reg}}q_{\mathrm{reg}}
+g^{(Q)}_{\mathrm{mob}}q_{\mathrm{mob}}
+g^{(Q)}_{\mathrm{dep}}q_{\mathrm{dep}},
\]
где \(q_{\mathrm{reg}}, q_{\mathrm{mob}}, q_{\mathrm{dep}}\) --- локальные уровни трёх режимов.
\[
z_{\mathrm{reg}}=v_1\widehat{Q}+v_2\widehat{V}+v_3\widehat{A}-v_4\widehat{\Gamma}_+,
\]
\[
z_{\mathrm{mob}}=v_5\widehat{Q}+v_6\widehat{V}-v_7\widehat{\Gamma}_+-v_8\widehat{A},
\]
\[
z_{\mathrm{dep}}=-v_9\widehat{Q}-v_{10}\widehat{V}-v_{11}\widehat{A}+v_{12}\widehat{\Gamma}_+.
\]
\[
\bigl(g_{\mathrm{reg}},g_{\mathrm{mob}},g_{\mathrm{dep}}\bigr)
=\softmax\!\left(\tau_I z_{\mathrm{reg}},\tau_I z_{\mathrm{mob}},\tau_I z_{\mathrm{dep}}\right).
\]
\[
R_6=g_{\mathrm{reg}}i_{\mathrm{reg}}+g_{\mathrm{mob}}i_{\mathrm{mob}}+g_{\mathrm{dep}}i_{\mathrm{dep}}.
\]
\[
H(g)=-
\frac{
g_{\mathrm{reg}}\ln g_{\mathrm{reg}}+
g_{\mathrm{mob}}\ln g_{\mathrm{mob}}+
g_{\mathrm{dep}}\ln g_{\mathrm{dep}}
}{\ln 3}.
\]
\[
B_6=g_{\mathrm{reg}}-g_{\mathrm{dep}}.
\]
\[
IIS_6=\clip\!\left(R_6-\lambda_H H(g)+\lambda_B B_6-\lambda_C C_{QV},0,1\right).
\]
Смысл версии: не один скалярный режим, а мягкая смесь трёх режимов: regulation, mobilization, depletion. \\

\bottomrule
\end{longtable}
\end{landscape}

\section{Практические модели}
\begin{landscape}
\small
\begin{longtable}{>{\raggedright\arraybackslash}p{1.6cm} >{\raggedright\arraybackslash}p{16.7cm}}
\toprule
\textbf{Версия} & \textbf{Точная практическая формула из кода} \\
\midrule
\endfirsthead
\toprule
\textbf{Версия} & \textbf{Точная практическая формула из кода} \\
\midrule
\endhead

\textbf{V3} &
\[
A=\frac{N_L-N_R}{N_L+N_R+\varepsilon},
\qquad
\Gamma=\ln(1+P_\gamma\cdot 10^{12}),
\qquad
V=\ln\!\left(1+\frac{HF}{LF}\right).
\]
\[
Q=\tanh(2.0\,Q_{\mathrm{raw}}),
\qquad
Q_{\mathrm{raw}}=\frac{\alpha_L-\alpha_R}{\alpha_L+\alpha_R+\varepsilon}.
\]
\[
K=\frac{R\cdot G}{1+1.2S^2+0.24F(S)},
\qquad
F(S)=\frac{1}{1+e^{-5.0(S-0.6)}}.
\]
\[
H=\frac{K\cdot DA}{50}-\frac{C}{15},
\qquad
IIS_3=\sigma\!\left(0.35A+0.25\Gamma+0.30H+0.10V\right).
\]
\textbf{Практическое ограничение:} в строгом режиме на открытых датасетах версия не является полностью direct-моделью, потому что \(DA\) и \(C\) напрямую не измеряются. \\

\textbf{V4} &
\[
A=\tanh\!\left(
1.10\,\Tnorm{\alpha_{\mathrm{asym}}}{0.0006126059}{0.0005296437}
0.40\,\Tnorm{\dfrac{P_L-P_R}{P_L+P_R+\varepsilon}}{-0.0298851838}{0.0287095439}
\right).
\]
\[
\Gamma=\tanh\!\left(
0.30\,\Tnorm{\ln(1+P_\gamma\cdot 10^{12})}{1.5434688676}{0.8496713487}
0.70\,\Tnorm{\ln(1+\gamma/\alpha)}{0.0003674267}{0.0004343448}
\right).
\]
\[
V=\tanh\!\left(
0.80\,\Tnorm{118.7401955561-HR}{0}{46.7184548147}
0.20\,\Tnorm{\ln(1+HF/LF)}{0.9675457417}{1.2220370049}
\right).
\]
\[
Q_4=\tanh\!\left(0.55A+0.45V-0.25\Gamma\right).
\]
\[
IIS_4=\sigma\!\left(0.10A-0.05\Gamma+0.25V+0.60Q_4\right).
\]
\textbf{Практическое ограничение:} если доступна только ЧСС или только \(HF/LF\), блок \(V\) упрощается до доступной части. \\

\textbf{V5} &
\[
\widehat{A}=\spow{A}{0.82},\qquad
\widehat{V}=\spow{V}{0.82},\qquad
\widehat{\Gamma}=\spow{\Gamma}{1.05}.
\]
\[
Q_{\mathrm{lin}}=0.50A+0.42V-0.30\Gamma,
\qquad
Q_{\mathrm{coh}}=AV-0.50|A-V|-0.18\Gamma_+.
\]
\[
S_{AV}=\sgn(\widehat{A}+\widehat{V})\sqrt{|\widehat{A}\widehat{V}|+\varepsilon},
\qquad
E_{AV}=\frac{|\widehat{A}|+|\widehat{V}|}{2},
\]
\[
C_{AV}=|\widehat{A}-\widehat{V}|+0.16\widehat{\Gamma}_+(1-\widehat{V}_+).
\]
\[
Q_5=\tanh\!\left(
1.4Q_{\mathrm{lin}}+0.55Q_{\mathrm{coh}}+0.42S_{AV}+0.28E_{AV}\sgn(Q_{\mathrm{lin}})-0.24C_{AV}
\right).
\]
\[
\widetilde{A}=\spow{A}{0.84},\qquad
\widetilde{V}=\spow{V}{0.84},\qquad
\widetilde{Q}=\spow{Q_5}{0.84}.
\]
\[
z_5=0.14\widetilde{A}-0.08\Gamma+0.22\widetilde{V}+0.56\widetilde{Q}-0.08.
\]
\[
S_{\mathrm{reg}}=
\sgn(\widetilde{A}+\widetilde{V})\sqrt{|\widetilde{A}\widetilde{V}|+\varepsilon}
+0.8\,\sgn(\widetilde{Q}+\widetilde{V})\sqrt{|\widetilde{Q}\widetilde{V}|+\varepsilon}.
\]
\[
E_{\mathrm{phase}}=\sqrt{\frac{\widetilde{A}^2+\widetilde{V}^2+\widetilde{Q}^2}{3}},
\qquad
C_{\mathrm{phase}}=|\widetilde{A}-\widetilde{V}|+0.75|\widetilde{Q}-\widetilde{V}|.
\]
\[
contrast_5=
2.6z_5+0.8z_5^3+0.18(Q_5-V)-0.06\Gamma_+
+0.52S_{\mathrm{reg}}+0.34E_{\mathrm{phase}}\sgn(z_5)-0.22C_{\mathrm{phase}}-0.10\Gamma_+(1-V_+).
\]
\[
IIS_{5,\mathrm{base}}=\sigma\!\left(core_5+0.22S_{\mathrm{reg}}-0.10C_{\mathrm{phase}}\right),
\qquad
core_5=0.12A-0.07\Gamma+0.24V+0.57Q_5.
\]
\[
IIS_{5,\mathrm{contrast}}=0.5+0.5\tanh(contrast_5).
\]
\[
IIS_5=\clip\!\left(0.65IIS_{5,\mathrm{base}}+0.35IIS_{5,\mathrm{contrast}},0,1\right).
\]
\textbf{Практическая цель:} убрать плоскость шкалы и усилить геометрию режимов без гормонов. \\

\textbf{V6} &
\[
\widehat{A}=\spow{A}{0.78},\qquad
\widehat{V}=\spow{V}{0.78},\qquad
\widehat{\Gamma}=\spow{\Gamma}{1.05},\qquad
\widehat{Q}=\spow{Q_6}{0.78}.
\]
\[
z^{(Q)}_{\mathrm{reg}}=0.95\widehat{A}+1.05\widehat{V}-0.55\widehat{\Gamma}_+,
\]
\[
z^{(Q)}_{\mathrm{mob}}=-0.18\widehat{A}+0.48\widehat{V}-1.05\widehat{\Gamma}_+,
\]
\[
z^{(Q)}_{\mathrm{dep}}=-0.80\widehat{A}-0.95\widehat{V}+0.72\widehat{\Gamma}_+.
\]
\[
\bigl(g^{(Q)}_{\mathrm{reg}},g^{(Q)}_{\mathrm{mob}},g^{(Q)}_{\mathrm{dep}}\bigr)
=\softmax\!\left(
1.65 z^{(Q)}_{\mathrm{reg}},
1.65 z^{(Q)}_{\mathrm{mob}},
1.65 z^{(Q)}_{\mathrm{dep}}
\right).
\]
\[
S_{AV}^{(6)}=\sgn(\widehat{A}+\widehat{V})\sqrt{|\widehat{A}\widehat{V}|+\varepsilon}.
\]
\[
q_{\mathrm{reg}}=0.76+0.18\tanh\!\left(1.10\widehat{A}+1.15\widehat{V}-0.45\widehat{\Gamma}_++0.30S_{AV}^{(6)}\right),
\]
\[
q_{\mathrm{mob}}=0.50+0.14\tanh\!\left(0.10\widehat{A}+0.55\widehat{V}-1.20\widehat{\Gamma}_+\right),
\]
\[
q_{\mathrm{dep}}=0.18+0.16\tanh\!\left(-0.65|\widehat{A}|-0.95|\widehat{V}|+0.95\widehat{\Gamma}_+\right).
\]
\[
Q_6=\clip\!\left(
g^{(Q)}_{\mathrm{reg}}q_{\mathrm{reg}}+
g^{(Q)}_{\mathrm{mob}}q_{\mathrm{mob}}+
g^{(Q)}_{\mathrm{dep}}q_{\mathrm{dep}},
0,1\right).
\]
\[
z_{\mathrm{reg}}=0.60\widehat{Q}+0.45\widehat{V}+0.15\widehat{A}-0.30\widehat{\Gamma}_+,
\]
\[
z_{\mathrm{mob}}=-0.05\widehat{Q}+0.20\widehat{V}-0.55\widehat{\Gamma}_+-0.08\widehat{A},
\]
\[
z_{\mathrm{dep}}=-0.70\widehat{Q}-0.55\widehat{V}-0.20\widehat{A}+0.35\widehat{\Gamma}_+.
\]
\[
\bigl(g_{\mathrm{reg}},g_{\mathrm{mob}},g_{\mathrm{dep}}\bigr)
=\softmax\!\left(1.65z_{\mathrm{reg}},1.65z_{\mathrm{mob}},1.65z_{\mathrm{dep}}\right).
\]
\[
S_{QV}^{(6)}=\sgn(\widehat{Q}+\widehat{V})\sqrt{|\widehat{Q}\widehat{V}|+\varepsilon},
\qquad
C_{QV}^{(6)}=|\widehat{A}-\widehat{V}|+0.75|\widehat{Q}-\widehat{V}|.
\]
\[
i_{\mathrm{reg}}=
0.80+0.16\tanh\!\left(
0.75\widehat{Q}+0.55\widehat{V}+0.18\widehat{A}-0.25\widehat{\Gamma}_++0.25S_{QV}^{(6)}
\right),
\]
\[
i_{\mathrm{mob}}=
0.54+0.12\tanh\!\left(
-0.08\widehat{Q}+0.32\widehat{V}-0.52\widehat{\Gamma}_++0.08\widehat{A}
\right),
\]
\[
i_{\mathrm{dep}}=
0.20+0.16\tanh\!\left(
-0.62\widehat{Q}-0.52\widehat{V}-0.18\widehat{A}+0.38\widehat{\Gamma}_+
\right).
\]
\[
R_6=g_{\mathrm{reg}}i_{\mathrm{reg}}+g_{\mathrm{mob}}i_{\mathrm{mob}}+g_{\mathrm{dep}}i_{\mathrm{dep}}.
\]
\[
H(g)=-
\frac{
g_{\mathrm{reg}}\ln g_{\mathrm{reg}}+
g_{\mathrm{mob}}\ln g_{\mathrm{mob}}+
g_{\mathrm{dep}}\ln g_{\mathrm{dep}}
}{\ln 3},
\qquad
B_6=g_{\mathrm{reg}}-g_{\mathrm{dep}}.
\]
\[
IIS_6=\clip\!\left(
R_6-0.08H(g)+0.05B_6-0.10C_{QV}^{(6)},
0,1\right).
\]
\[
raw_6=0.14\widehat{A}-0.08\widehat{\Gamma}+0.24\widehat{V}+0.54\widehat{Q}.
\]
\textbf{Практическая цель:} уйти от одной почти линейной формулы к мягкой смеси трёх режимов. \\

\bottomrule
\end{longtable}
\end{landscape}

\section{Практико-ориентированные (proxy) модели}
\begin{landscape}
\small
\begin{longtable}{>{\raggedright\arraybackslash}p{1.6cm} >{\raggedright\arraybackslash}p{16.7cm}}
\toprule
\textbf{Версия} & \textbf{Proxy-формула и её смысл} \\
\midrule
\endfirsthead
\toprule
\textbf{Версия} & \textbf{Proxy-формула и её смысл} \\
\midrule
\endhead

\textbf{V3} &
\[
R=
\begin{cases}
\dfrac{HF/LF}{1+HF/LF}, & \text{если доступно } HF/LF,\\[4pt]
\dfrac{HF}{1+HF}, & \text{если доступно только } HF,
\end{cases}
\]
\[
G=\frac{\ln(1+P_\gamma\cdot 10^{12})}{1+\ln(1+P_\gamma\cdot 10^{12})},
\qquad
S=\clip\!\left(\frac{C_{\mathrm{proxy}}-8}{25-8},0,1\right),
\]
\[
F(S)=\frac{1}{1+e^{-5.0(S-0.6)}},
\qquad
K=\frac{R\cdot G}{1+1.2S^2+0.24F(S)}.
\]
\[
H_{\mathrm{proxy}}=\frac{K\cdot DA_{\mathrm{proxy}}}{50}-\frac{C_{\mathrm{proxy}}}{15}.
\]
\[
IIS_{3,\mathrm{proxy}}=\sigma\!\left(0.35A+0.25\Gamma+0.30H_{\mathrm{proxy}}+0.10V\right).
\]
Здесь \(DA_{\mathrm{proxy}}\) и \(C_{\mathrm{proxy}}\) берутся из self-report, stress-меток и нормированных косвенных признаков. \\

\textbf{V4} &
\[
Q_{4,\mathrm{anchor}}=
\tanh\!\left(
2.2\left(
0.62\,valence_{\mathrm{norm}}+0.38(1-arousal_{\mathrm{norm}})-0.5
\right)\right).
\]
\[
Q_{4,\mathrm{proxy}}=0.65Q_4+0.35Q_{4,\mathrm{anchor}}.
\]
\[
IIS_{4,\mathrm{proxy}}=\sigma\!\left(0.10A-0.05\Gamma+0.25V+0.60Q_{4,\mathrm{proxy}}\right).
\]
Proxy здесь не заменяет физиологию целиком, а только мягко якорит \(Q\) к реальным valence/arousal. \\

\textbf{V5} &
\[
Q_{5,\mathrm{anchor}}=
\tanh\!\left(
2.4\left(
0.65\,valence_{\mathrm{norm}}+0.35(1-arousal_{\mathrm{norm}})-0.5
\right)\right).
\]
\[
Q_{5,\mathrm{proxy}}=0.75Q_5+0.25Q_{5,\mathrm{anchor}}.
\]
\[
IIS_{5,\mathrm{proxy}}=
\clip\!\left(
0.65IIS_{5,\mathrm{base}}(Q_{5,\mathrm{proxy}})
+0.35IIS_{5,\mathrm{contrast}}(Q_{5,\mathrm{proxy}}),
0,1\right).
\]
Proxy-вклад уменьшен по сравнению с V4, чтобы не разрушать собственную геометрию нелинейной модели. \\

\textbf{V6} &
\[
Q_{6,\mathrm{anchor}}=
0.70\,valence_{\mathrm{norm}}+0.30(1-arousal_{\mathrm{norm}}).
\]
\[
Q_{6,\mathrm{proxy}}=0.82Q_6+0.18Q_{6,\mathrm{anchor}}.
\]
\[
\widehat{Q}_{\mathrm{proxy}}=\spow{Q_{6,\mathrm{proxy}}}{0.78}.
\]
Дальше вся gated-структура версии 6 пересчитывается уже с \(Q_{6,\mathrm{proxy}}\):
\[
IIS_{6,\mathrm{proxy}}=
\clip\!\left(
R_6(Q_{6,\mathrm{proxy}})
-0.08H(g)
+0.05B_6
-0.10C_{QV}^{(6)}(Q_{6,\mathrm{proxy}}),
0,1\right).
\]
Proxy-режим V6 оставляет архитектуру режимов, но слегка подтягивает \(Q\) к внешним оценкам состояния. \\

\bottomrule
\end{longtable}
\end{landscape}

\section{Приложение: динамический слой поверх версий V4--V6}
Текущая динамика в проекте строится \textbf{поверх} статического IIS и не меняет базовую формулу версии. Для окна \(t\):
\[
\overline{I}_t=\operatorname{rolling\_mean}(IIS_{1:t}),
\qquad
\widetilde{I}_t=\operatorname{rolling\_median}(IIS_{1:t}),
\]
\[
I^{\mathrm{ewma}}_t=\operatorname{EWMA}_{\alpha=0.35}(IIS_{1:t}),
\qquad
I^{\mathrm{fast}}_t=\operatorname{EWMA}_{\alpha=0.42}(IIS_{1:t}).
\]
\[
vol_t=\operatorname{rolling\_mean}\bigl(|IIS_t-IIS_{t-1}|\bigr).
\]
\[
w_t=\clip\!\left(\frac{vol_t}{0.035},\,0.10,\,0.70\right).
\]
\[
core_t=0.60\,I^{\mathrm{fast}}_t+0.40\,\widetilde{I}_t.
\]
\[
target_t=(1-w_t)\,IIS_t+w_t\,core_t.
\]
\[
stress_t=0.45\frac{\Gamma_t+1}{2}+0.35\frac{1-V_t}{2}+0.20\frac{1-Q_t}{2}.
\]
\[
\alpha^{\mathrm{resp}}_t=\clip\!\left(0.48+0.22\,stress_t+0.10\,\frac{vol_t}{0.035},\,0.08,\,0.85\right),
\]
\[
\alpha^{\mathrm{rec}}_t=\clip\!\left(0.12+0.10(1-stress_t)-0.04\,\frac{vol_t}{0.035},\,0.08,\,0.85\right).
\]
\[
IIS^{\mathrm{dyn}}_t=
\begin{cases}
IIS_t, & t=1,\\[4pt]
IIS^{\mathrm{dyn}}_{t-1}+\alpha_t\bigl(target_t-IIS^{\mathrm{dyn}}_{t-1}\bigr), & t>1,
\end{cases}
\]
где
\[
\alpha_t=
\begin{cases}
\alpha^{\mathrm{resp}}_t, & target_t\le IIS^{\mathrm{dyn}}_{t-1},\\[4pt]
\alpha^{\mathrm{rec}}_t, & target_t> IIS^{\mathrm{dyn}}_{t-1}.
\end{cases}
\]
Это причинная схема: ухудшение проходит быстрее, восстановление медленнее.

\section{Итог}
В текущем проекте:
\begin{itemize}[leftmargin=1.2cm]
    \item \textbf{V3} --- последняя гормонально-зависимая версия с коэффициентом \(K\);
    \item \textbf{V4} --- базовая практическая direct-версия без гормонов;
    \item \textbf{V5} --- нелинейная контрастная версия без гормонов;
    \item \textbf{V6} --- универсальная gated-версия без гормонов.
\end{itemize}

Таким образом документ разделяет:
\[
\text{теорию модели}
\;\longrightarrow\;
\text{реальную direct-реализацию}
\;\longrightarrow\;
\text{proxy-реализацию при неполных данных}.
\]

\end{document}
